\chapter{Cryotherapy}

\section*{Overview}
Cryotherapy uses extreme cold (typically liquid nitrogen at -196 C) to destroy abnormal or diseased tissue, such as warts, skin tags, or precancerous skin lesions.

\section*{Alternative Names}
Cryosurgery, Skin freezing

\section*{Principle}
Extreme cold produced by liquid nitrogen, cryoprobes, or other cryogens is applied to abnormal tissue, causing intracellular ice formation, cell membrane damage, and osmotic injury. These effects kill the target cells while surrounding tissue heals by secondary intention.
\section*{History}
Dr. James Arnott first used cold mixtures to treat tumors in 1845. Liquid nitrogen became the standard cryogen in the 1940s-1950s.

\section*{Why It Is Done}
Ordered to treat common warts, seborrheic keratoses, skin tags, actinic keratoses (precancerous lesions), or superficial basal cell carcinomas.

\section*{Is This a Primary or Secondary Test or Procedure}
Primary treatment for common viral warts and actinic keratoses.

\section*{How It Is Performed}
Liquid nitrogen is applied to the target skin lesion using a cotton-tipped applicator or a pressurized spray gun. The lesion is frozen for several seconds, allowed to thaw, and sometimes frozen again.

\section*{Preparation}
No special preparation. The skin is cleaned before the procedure.

\section*{Contraindications and Precautions}
Lesions suspected of melanoma (biopsy must be done instead), or patients with cold urticaria or cryoglobulinemia.

\section*{Risks}
Local pain, blistering, hypopigmentation (white mark, common), hyperpigmentation, or scar formation.

\section*{Aftercare and Recovery}
A blister will form at the site, which will crust over and peel off in 1-2 weeks. Keep the area clean and do not pop the blister.

\section*{Outcome and Follow-up}
Successful treatment results in the sloughing off of the abnormal tissue and replacement with normal skin.

\section*{Expected Results}
Destruction of the target lesion; healing of the skin with minimal scarring or discoloration.

\section*{Follow-up}
Re-evaluation in 2-3 weeks; repeat freezing session if the lesion is not fully resolved.

\section*{When to Seek Medical Care}
Follow the procedure team's specific instructions. Seek urgent medical assessment for severe or worsening pain, uncontrolled bleeding, fever or signs of infection, trouble breathing, chest pain, fainting, new weakness or confusion, inability to urinate, or any rapidly worsening symptom. The appropriate warning signs depend on the procedure and the patient's medical condition.

\section*{References}
\begin{enumerate}
  \item MedlinePlus. Cryotherapy. U.S. National Library of Medicine; 2025.
  \item Current Medical Diagnosis and Treatment. McGraw-Hill; 2025.
\end{enumerate}

\clearpage
